
\documentclass{article}
\usepackage{colortbl}
\usepackage{makecell}
\usepackage{multirow}
\usepackage{supertabular}

\begin{document}

\newcounter{utterance}

\centering \large Interaction Transcript for game `eqbench', experiment `eq\_bench', episode 43 with reflexion{-}Qwen3.5{-}27B.
\vspace{24pt}

{ \footnotesize  \setcounter{utterance}{1}
\setlength{\tabcolsep}{0pt}
\begin{supertabular}{c@{$\;$}|p{.15\linewidth}@{}p{.15\linewidth}p{.15\linewidth}p{.15\linewidth}p{.15\linewidth}p{.15\linewidth}}
    \# & \multicolumn{2}{c}{Player} && \multicolumn{2}{c}{Game Master} \\
    \hline

    \theutterance \stepcounter{utterance}  
    & & & \multicolumn{4}{p{0.6\linewidth}}{
        \cellcolor[rgb]{0.9,0.9,0.9}{
            \makecell[{{p{\linewidth}}}]{
                \texttt{\tiny{[P1$\langle$GM]}}
                \texttt{You have attempted this same task before and did not succeed. These are the plans you wrote after those attempts.} \\
\texttt{Trial \#1: New plan: My previous attempt failed because I assigned a high score to "Determined" while overlooking the strong textual evidence for "Sympathetic." Rachel's dialogue explicitly shows empathy ("It's not like I do it on purpose," "We're going to have to figure it out") and she actively pushes for a solution to help James, which is the definition of determination in this context. However, her tone is supportive, not just goal{-}oriented. I should have recognized that her primary emotional state is a blend of sympathy for his struggle and determination to resolve it for the play. In the future, I will re{-}read the final lines specifically to identify the character's intent toward the other person. If the character offers help or reassurance, I will prioritize "Sympathetic" and "Determined" as the top scores, ensuring the highest score reflects the most dominant emotional driver (likely Sympathetic here, given the "intimidating" and "retreating" context), rather than defaulting to a generic "Determined" score. I will also ensure "Dismissive" remains low as she engages directly, and "Nervous" is low as she appears confident.} \\
\\ 
\texttt{Trial \#2: New plan: My previous attempt failed because I overestimated "Determined" (7) and underestimated the nuance of "Sympathetic" (8) relative to the specific context of Rachel's final lines. While she is determined to rehearse, her primary emotional driver in the final exchange ("We're going to have to figure it out, aren't we?") is a supportive, problem{-}solving empathy toward James's vulnerability, not just a rigid goal orientation. I incorrectly treated "Determined" as a high{-}intensity emotion comparable to her sympathy, whereas the text suggests her determination is a *method* of expressing sympathy. In the future, I will distinguish between the *goal* (rehearsing) and the *emotional tone* (supportive empathy). I will assign the highest score to "Sympathetic" (e.g., 9) as it encompasses her understanding of his "haunted" state and her offer to help. I will lower "Determined" (e.g., to 4 or 5) to reflect that it is a secondary action rather than the primary emotional state. I will keep "Nervous" and "Dismissive" at 0 or 1, as her confident, direct engagement contradicts these emotions. The key correction is to ensure the highest score reflects the character's *attitude toward the other person* (sympathy) rather than just their *intent to act* (determination).} \\
\\ 
\\ 
\\ 
\texttt{Your task is to predict the likely emotional responses of a character in this dialogue:} \\
\\ 
\texttt{James: I can't believe we're doing this again, Rachel. } \\
\texttt{Rachel: What, the play? Or you mean us, playing lovers?} \\
\texttt{James: Both, I guess. } \\
\texttt{Rachel: Well, it's not like we have a choice. We were cast for these roles. } \\
\texttt{James: Yeah, but it's... complicated.} \\
\texttt{Rachel: Complicated? Are you still stuck on that scene from three years ago?} \\
\texttt{James: I wouldn't say stuck. More like... haunted.} \\
\texttt{Rachel: Haunted? That's a bit dramatic, even for you, James.} \\
\texttt{James: Maybe, but it's the truth.} \\
\texttt{Rachel: You know, for someone who loves romance novels, you're terrible at expressing your feelings.} \\
\texttt{James: And you're too good at it. It's intimidating.} \\
\texttt{Rachel: Intimidating? Me? You're the one who's always retreating into your shell.} \\
\texttt{James: It's not like I do it on purpose, Rachel. It's just... hard.} \\
\texttt{Rachel: What's hard, James? Acting or being around me?} \\
\texttt{James: Both.} \\
\texttt{Rachel: Well, we're going to have to figure it out, aren't we? For the play.} \\
\texttt{James: Yes, for the play.} \\
\texttt{Rachel: So, let's rehearse the scene. The kiss scene.} \\
\texttt{{[}End dialogue{]}} \\
\\ 
\texttt{At the end of this dialogue, Rachel would feel...} \\
\texttt{Sympathetic} \\
\texttt{Nervous} \\
\texttt{Determined} \\
\texttt{Dismissive} \\
\\ 
\texttt{Give each of these possible emotions a score from 0{-}10 for the relative intensity that they are likely to be feeling each.} \\
\\ 
\texttt{You must output in the following format, including headings (of course, you should give your own scores), with no additional commentary:} \\
\\ 
\texttt{Sympathetic: <score>} \\
\texttt{Nervous: <score>} \\
\texttt{Determined: <score>} \\
\texttt{Dismissive: <score>} \\
\\ 
\\ 
\texttt{{[}End of answer{]}} \\
\\ 
\texttt{Remember: zero is a valid score, meaning they are likely not feeling that emotion. You must score at least one emotion > 0.} \\
\\ 
\texttt{Your answer:} \\
            }
        }
    }
     \\ \\

    \theutterance \stepcounter{utterance}  
    & \multicolumn{4}{p{0.6\linewidth}}{
        \cellcolor[rgb]{0.561,0.737,0.561}{
            \makecell[{{p{\linewidth}}}]{
                \texttt{\tiny{[P1$\rangle$GM]}}
                \texttt{Sympathetic: 9} \\
\texttt{Nervous: 0} \\
\texttt{Determined: 6} \\
\texttt{Dismissive: 1} \\
            }
        }
    }
    & & \\ \\

    \theutterance \stepcounter{utterance}  
    & & & \multicolumn{2}{p{0.3\linewidth}}{
        \cellcolor[rgb]{0.9,0.9,0.9}{
            \makecell[{{p{\linewidth}}}]{
                \texttt{\tiny{[GM$|$GM]}}
                \texttt{Sympathetic: 9} \\
\texttt{Nervous: 0} \\
\texttt{Determined: 6} \\
\texttt{Dismissive: 1} \\
            }
        }
    }
    & & \\ \\

    \theutterance \stepcounter{utterance}  
    & & & \multicolumn{2}{p{0.3\linewidth}}{
        \cellcolor[rgb]{0.9,0.9,0.9}{
            \makecell[{{p{\linewidth}}}]{
                \texttt{\tiny{[GM$|$GM]}}
                \texttt{Sympathetic: 4} \\
\texttt{Nervous: 0} \\
\texttt{Determined: 8} \\
\texttt{Dismissive: 2} \\
            }
        }
    }
    & & \\ \\

    \theutterance \stepcounter{utterance}  
    & & & \multicolumn{2}{p{0.3\linewidth}}{
        \cellcolor[rgb]{0.9,0.9,0.9}{
            \makecell[{{p{\linewidth}}}]{
                \texttt{\tiny{[GM$|$GM]}}
                \texttt{game\_result = LOSE} \\
            }
        }
    }
    & & \\ \\

    \theutterance \stepcounter{utterance}  
    & & \multicolumn{4}{p{0.6\linewidth}}{
        \cellcolor[rgb]{0.95,0.95,0.95}{
            \makecell[{{p{\linewidth}}}]{
                \texttt{\tiny{[P1$\leftrightarrow$P]}}
                \texttt{Post{-}episode Reflexion (not sent to GM):} \\
\texttt{New plan: My previous attempts failed because I misinterpreted the character's pragmatic push to "rehearse" as the primary emotion "Determined," assigning it a score (6) that was too high relative to the context. I overlooked that Rachel's determination is a *tool* she uses to express her underlying "Sympathetic" nature toward James's vulnerability. The text shows her actively managing his anxiety ("It's not like I do it on purpose," "We're going to have to figure it out"), which is the definition of sympathy, not just goal{-}oriented determination. In the future, I will distinguish between the *action* (rehearsing) and the *emotional driver* (empathy). I will assign the highest score to "Sympathetic" (e.g., 9 or 10) as it is the dominant tone. I will significantly lower "Determined" (e.g., to 2 or 3) to reflect that it is a secondary, functional response rather than the core emotional state. I will keep "Nervous" and "Dismissive" at 0, as her confident, direct engagement contradicts these. The key correction is to ensure the highest score reflects the character's *attitude toward the other person's feelings* rather than their *intent to complete a task*.} \\
            }
        }
    }
    & \\ \\

\end{supertabular}
}

\end{document}
